\documentclass[12pt]{report}
\usepackage[utf8]{inputenc}
\usepackage[russian]{babel}
%\usepackage[14pt]{extsizes}
\usepackage{listings}
\usepackage{graphicx}
\usepackage{amsmath,amsfonts,amssymb,amsthm,mathtools} 
\usepackage{pgfplots}
\usepackage{filecontents}
\usepackage{float}
\usepackage{indentfirst}
\usepackage{eucal}
\usepackage{enumitem}
%s\documentclass[openany]{book}
\frenchspacing

\usepackage{indentfirst} % Красная строка

\usetikzlibrary{datavisualization}
\usetikzlibrary{datavisualization.formats.functions}

\usepackage{amsmath}


% Для листинга кода:
\lstset{ %
	language=c,                 % выбор языка для подсветки (здесь это С)
	basicstyle=\small\sffamily, % размер и начертание шрифта для подсветки кода
	numbers=left,               % где поставить нумерацию строк (слева\справа)
	numberstyle=\tiny,           % размер шрифта для номеров строк
	stepnumber=1,                   % размер шага между двумя номерами строк
	numbersep=5pt,                % как далеко отстоят номера строк от подсвечиваемого кода
	showspaces=false,            % показывать или нет пробелы специальными отступами
	showstringspaces=false,      % показывать или нет пробелы в строках
	showtabs=false,             % показывать или нет табуляцию в строках
	frame=single,              % рисовать рамку вокруг кода
	tabsize=2,                 % размер табуляции по умолчанию равен 2 пробелам
	captionpos=t,              % позиция заголовка вверху [t] или внизу [b] 
	breaklines=true,           % автоматически переносить строки (да\нет)
	breakatwhitespace=false, % переносить строки только если есть пробел
	escapeinside={\#*}{*)}   % если нужно добавить комментарии в коде
}


\usepackage[left=2cm,right=2cm, top=2cm,bottom=2cm,bindingoffset=0cm]{geometry}
% Для измененных титулов глав:
\usepackage{titlesec, blindtext, color} % подключаем нужные пакеты
\definecolor{gray75}{gray}{0.75} % определяем цвет
\newcommand{\hsp}{\hspace{20pt}} % длина линии в 20pt
% titleformat определяет стиль
\titleformat{\chapter}[hang]{\Huge\bfseries}{\thechapter\hsp\textcolor{gray75}{|}\hsp}{0pt}{\Huge\bfseries}


% plot
\usepackage{pgfplots}
\usepackage{filecontents}
\usetikzlibrary{datavisualization}
\usetikzlibrary{datavisualization.formats.functions}

\begin{document}
	%\def\chaptername{} % убирает "Глава"
	\thispagestyle{empty}
	\begin{titlepage}
		\noindent \begin{minipage}{0.15\textwidth}
			\includegraphics[width=\linewidth]{b_logo}
		\end{minipage}
		\noindent\begin{minipage}{0.9\textwidth}\centering
			\textbf{Министерство науки и высшего образования Российской Федерации}\\
			\textbf{Федеральное государственное бюджетное образовательное учреждение высшего образования}\\
			\textbf{~~~«Московский государственный технический университет имени Н.Э.~Баумана}\\
			\textbf{(национальный исследовательский университет)»}\\
			\textbf{(МГТУ им. Н.Э.~Баумана)}
		\end{minipage}
		
		\noindent\rule{18cm}{3pt}
		\newline\newline
		\noindent ФАКУЛЬТЕТ $\underline{\text{«Информатика и системы управления»}}$ \newline\newline
		\noindent КАФЕДРА $\underline{\text{«Программное обеспечение ЭВМ и информационные технологии»}}$\newline\newline\newline\newline\newline
		
		\begin{center}
			\noindent\begin{minipage}{1.1\textwidth}\centering
				\Large\textbf{  Отчет по лабораторной работе №1}\newline
				\textbf{по дисциплине <<Функциональное и логическое}\newline
				\textbf{~~~программирование>>}\newline\newline
			\end{minipage}
		\end{center}
		
		\noindent\textbf{Тема} $\underline{\text{Списки в  Lisp. Использование стандартных функций~~}}$\newline\newline
		\noindent\textbf{Студент} $\underline{\text{Варламова Е.А.~~~~~~~~~~~~~~~}}$\newline\newline
		\noindent\textbf{Группа} $\underline{\text{ИУ7-61Б~~~~~~~~~~~~~~~~~~~~~~~}}$\newline\newline
		\noindent\textbf{Оценка (баллы)} $\underline{\text{~~~~~~~~~~~~~~~~~~~~~~}}$\newline\newline
		\noindent\textbf{Преподаватель} $\underline{\text{Толпинская Н.Б.}}$\newline\newline\newline
		
		\begin{center}
			\vfill
			Москва~---~\the\year
			~г.
		\end{center}
	\end{titlepage}
\setcounter{page}{2}

\section*{Задание 1}
Представить следующие списки в виде списочных ячеек:
	\begin{enumerate}
		\item ’(open close halph)
		
		\item ’((open1) (close2) (halph3))
		
		\item ’((one) for all (and (me (for you))))
		
		\item ’((TOOL)(call))
		
		\item ’((TOOL1)((call2))((sell)))
		
		\item ’(((TOOL)(call))(sell))
	\end{enumerate}
	
	\textbf{Решение}
	Решение оформлено на тетрадном листе бумаге, прилагающемуся к отчету.	
	
\section*{Задание 2}

Используя только функции CAR и CDR, написать выражения, возвращающие второй, третий, четвертый элементы заданного списка.
	
\textbf{Решение}

\begin{lstlisting}[label=first,caption=Решение задания №2, language=lisp]
(car (cdr '(1 2 3 4)))
(car (cdr (cdr '(1 2 3 4))))
(car (cdr (cdr (cdr '(1 2 3 4)))))
\end{lstlisting}

\section*{Задание №3}

Что будет в результате вычисления выражений?

\textbf{Решение}

\begin{lstlisting}[label=second,caption=Решение задания №3, language=lisp]
(caadr '((blue cube) (red pyaramid))) ; red
(cdar '((abc) (def) (ghi))) ; Nil
(cadr '((abc) (def) (ghi))) ; (def)
(caddr '((abc) (def) (ghi))) ; (ghi)
\end{lstlisting}

\section*{Задание 4}
Напишите результат вычисления выражений:

\textbf{Решение}

\begin{lstlisting}[label=second,caption=Решение задания №4, language=lisp]
(list 'Fred 'and  'Wilma) ; (Fred and Wilma)
(list 'Fred '(and Wilma)) ; (Fred (and Wilma))
(cons Nil Nil) ; (Nil)
(cons T Nil) ; (T)
(cons Nil T) ; (Nil . T)
(list Nil) ; (Nil)
(cons '(T) Nil) ; ((T))
(list '(one two) '(free temp)) ; ((one two) (free temp))

(cons 'Fred '(and Wilma)) ; (Fred and Wilma)
(cons 'Fred '(Wilma)) ; (Fred Wilma)
(list Nil Nil) ; (Nil Nil)
(list T Nil) ; (T Nil)
(list Nil T) ; (Nil T)
(cons T (list Nil)) ; (T Nil)
(list '(T) Nil) ; ((T) Nil)
(cons '(one two) '(free temp)) ; ((one two) free temp)
\end{lstlisting}

\section*{Задание №5}

\indent Написать функцию (f ar1 ar2 ar3 ar4), возвращаущую список: ((ar1 ar) (ar3 ar4)).""\newline
\indent Написать функцию (f ar1 ar2), возвращаущую ((ar1) (ar2)).""\newline
\indent Написать функцию (f ar1), возвращаущую (((ar1))).""\newline
\indent Представить результаты в виде списочных ячеек. 
\textbf{Решение}
\begin{lstlisting}[label=third,caption=Решение задания №5 (функция №1), language=lisp]
(DEFUN f (ar1 ar2 ar3 ar4) (list (list ar1 ar2) (list ar3 ar4) ))
# (DEFUN f (ar1 ar2 ar3 ar4) (cons (cons ar1 (cons ar2 nil)) (cons (cons ar3 (cons ar4 nil)) nil) ))
\end{lstlisting}

\begin{lstlisting}[label=third,caption=Решение задания №5 (функция №2), language=lisp]
(defun f (ar1 ar2) (list (list ar1) (list ar2)) )
\end{lstlisting}

\begin{lstlisting}[label=third,caption=Решение задания №5 (функция №3), language=lisp]
(defun f (ar1) (list (list (list ar1))) )
\end{lstlisting}
Результаты в виде списочных ячеек приведены а прилагающихся к отчёту тетрадных листах.
\section*{Контрольные вопросы}
	\textbf{Вопрос 1.} Элементы языка: определение, синтаксис, представление в памяти.
	
	Элементами языка Lisp являются атомы и структуры (точечные пары, списки). К атомам относятся:
	\begin{itemize}
		\item символы -- набор литер, начинающихся с буквы.
		\item специальные символы: $\{T, Nil\}$ (используются для обозначения логических констант).
		\item самоопределимые атомы -- натуральные, дробные, вещественные числа, строки (последовательность символов, заключенных в двойные апострофы)
	\end{itemize}
	
	\newline
	
	Точечные пары ::= (<атом>, <атом>) | (<атом>, <точечная пара>) |
	
	(<точечная пара>, <атом>) | (<точечная пара>, <точечная пара>)
	
	""\newline
	\indent Список ::= <пустой список> | <непустой список>, где
	
	<пустой список> ::= () | Nil, <непустой список> ::= (<первый элемент>, <хвост>),
	
	<первый элемент> ::= <S-выражение>, <хвост> ::= <список>""\newline
	
	\indent \textbf{Список} -- частный случай S-выражения. Любая структура (точечная пара или список) заключаются в круглые скобки:
	
\begin{itemize}
	\item (A . B) -- точечная пара;
	\item $(A)$ -- список из одного элемента;
	\item $Nil$ или $()$ -- пустой список;
	\item (A . (B . (C . (D ()))))) или (A B C D) -- непустой список;
	\item Элементы списка могут являться списками: $((A)(B)(CD))$
\end{itemize}
	
Любая непустая структура в Lisp, в памяти представленна списковой ячейкой, хранящей два указателя: на голову и хвост.

	
\textbf{Вопрос 2.} Особенности языка Lisp. Структура программы. Символ апостроф.

Важной особенностью языка Lisp является единая синтаксическая форма записи программ и данных, что позволяет обрабатывать структуры данных как программы и модифицировать программы как данные. 

Lisp-программа представляет собой последовательность вычислимых выражений, явяляющихся атомом или списком.

Символ ' эквивалентен функции quote – он блокирует вычисление выражения. Таким образом, выражение воспринимается интерпретатором как данные.""\newline

\textbf{Вопрос 3.} Базис языка Lisp. Ядро языка. \newline

Базис языка образуют атомы, структуры (точечные пары и списки), базовые функции, базовые функционалы (функции, аргументами и значением которых являются функции).

Ядро Лиспа работает следующим образом: 
\begin{enumerate}
\item ожидает ввода S-выражения; 
\item передает введенное S-выражение функции EVAL (которая вычисляет значение своего единственного аргумента и возвращает его в качестве результата); 
\item выводит полученный результат;
\item переходит к пункту 1.
\end{enumerate}
	
\end{document}
\section*{Цель и задачи работы}
\textbf{Целью} работы является приобретение навыков использования списков и стандартных функций Lisp. Для достижения поставленной цели необходимо решить следующие \textbf{задачи}:
\begin{itemize}
    \item изучить способ использования списков для фиксации информации;
    \item изучить внутреннее представление одноуровневых и структурированных спсиков и методы их обработки с использованием базовых функций Lisp.
\end{itemize}